% !TEX root = ../../hw05.tex
\begin{proof}[Solution]
  Let
  \[
    p(x)=a_mx^m+a_{m-1}x^{m-1}+\cdots+a_0,
  \]
  where $m$ is odd and $a_m\ne0$. Dividing by $a_m$ does not change
  the roots, so we work with
  \[
    q(x)=\frac{p(x)}{a_m}
      =x^m+b_{m-1}x^{m-1}+\cdots+b_0.
  \]
  We want to choose two inputs where the leading term is larger in
  absolute value than all the remaining terms together.

  Set $C=\sum_{k=0}^{m-1}|b_k|$. For $|x|\ge1$, every power
  $|x|^k$ with $k\le m-1$ is at most $|x|^{m-1}$. The triangle
  inequality therefore gives
  \[
    \left|\sum_{k=0}^{m-1}b_kx^k\right|
      \le\sum_{k=0}^{m-1}|b_k||x|^k
      \le C|x|^{m-1}.
  \]
  Choose $R>\max\{1,C\}$. Since $m$ is odd, $(-R)^m=-R^m$.
  Applying the bound at $R$ and $-R$, we obtain
  \begin{align*}
    q(R)&\ge R^m-CR^{m-1}
      =R^{m-1}(R-C)>0,\\
    q(-R)&\le -R^m+CR^{m-1}
      =-R^{m-1}(R-C)<0.
  \end{align*}
  The polynomial $q$ is continuous on $[-R,R]$. By the intermediate
  value theorem, there is a $c\in(-R,R)$ such that $q(c)=0$.
  Since $p(c)=a_mq(c)$, we also have $p(c)=0$.
\end{proof}
